\lysh{22564_SOLUTION}{1}
\textbf{ΛΥΣΗ}

\alpha) Οι συντεταγμένες του σημείου $A(2,2)$ επαληθεύουν τις εξισώσεις των δύο ελλείψεων αφού
\[
\frac{2^2}{12} + \frac{2^2}{6} = \frac{4}{12} + \frac{4}{6} = \frac{1}{3} + \frac{2}{3} = 1
\]
και
\[
\frac{2^2}{6} + \frac{2^2}{12} = \frac{4}{6} + \frac{4}{12} = \frac{2}{3} + \frac{1}{3} = 1
\]
Το σημείο $B(-2,2)$ είναι συμμετρικό του σημείου $A(2,2)$ ως προς τον άξονα $y'y$, επομένως θα ανήκει στις δύο ελλείψεις, αφού και το σημείο $A(2,2)$ ανήκει στις δύο ελλείψεις.

% TODO: TikZ σχήμα

\beta) Η εφαπτομένη $\varepsilon_1$ της έλλειψης $C_1: \frac{x^2}{12} + \frac{y^2}{6} = 1$ στο σημείο $A(2,2)$ έχει εξίσωση:
\[
\frac{2x}{12} + \frac{2y}{6} = 1 \Leftrightarrow \frac{x}{6} + \frac{y}{3} = 1 \Leftrightarrow x + 2y = 6 \Leftrightarrow x + 2y - 6 = 0.
\]
Η εφαπτομένη $\varepsilon_2$ της έλλειψης $C_2: \frac{x^2}{6} + \frac{y^2}{12} = 1$ στο σημείο $B(-2,2)$ έχει εξίσωση:
\[
\frac{-2x}{6} + \frac{2y}{12} = 1 \Leftrightarrow \frac{-x}{3} + \frac{y}{6} = 1 \Leftrightarrow -2x + y = 6 \Leftrightarrow -2x + y - 6 = 0.
\]

\gamma) Ο συντελεστής διεύθυνσης της εφαπτομένης $\varepsilon_1: x+2y-6=0$ είναι $\lambda_1 = -\frac{1}{2}$ και ο συντελεστής διεύθυνσης της εφαπτομένης $\varepsilon_2: -2x+y-6=0$ είναι $\lambda_2 = 2$.
Οι εφαπτομένες $\varepsilon_1$, $\varepsilon_2$ είναι κάθετες, γιατί $\lambda_1 \lambda_2 = -\frac{1}{2} \cdot 2 = -1$.
}